% a06_WN_c: part 3 of 4 of the walled reading of Ba3 Prop. 1 (input from sec/a06_BALtorus_d7b.tex): the proof of thm:a06:WN(iii), the phi phi A node clauses with walled lines.

\begin{proof}[Proof of \cref{thm:a06:WN}\textup{(iii)}]
\emph{The shapes.} A dangerous $\phi\phi A$ node has order $3$ and one loop (Step~3(D$_4$)), and a node containing a finite \cite[(1.7)]{Ba3} vertex is not dangerous (\cref{lem:a06d:S25}(b)); by \cite[p.~429]{Ba3} its vertices are of the types \cite[(1.6)--(1.8), (1.10)]{Ba3} or divergent-order counterterm vertices. With an odd number of vector legs the vertex multisets are three one-vector-leg vertices \cite[(1.8)]{Ba3}; one one-vector-leg vertex with one two-vector-leg vertex of type \cite[(1.8)]{Ba3} or \cite[(1.10)]{Ba3}; one three-vector-leg vertex \cite[(1.8)]{Ba3}; or one one-vector-leg vertex with a \cite[(1.6)]{Ba3} vertex or with an order-$2$ counterterm vertex. The line count gives one loop in each case, so the internal lines form one cycle through all vertices (a triangle or a double line) in the bridgeless case, and a bridge otherwise. The single three-vector-leg vertex is excluded (a vector self-loop leaves a differentiated scalar slot, a scalar self-loop leaves three vector slots). The bridgeless shapes are: \emph{T}, three one-vector-leg vertices in a triangle, the apex carrying the vector slot and the two base vertices the scalar slots, which are their undifferentiated legs, the base vertices' differences acting on their lines to the apex, and the apex's difference on one of its two lines; \emph{B8}, resp.\ \emph{B10}, a one-vector-leg vertex $u$ and a two-vector-leg vertex $w$ of type \cite[(1.8)]{Ba3}, resp.\ \cite[(1.10)]{Ba3}, joined by a vector line and a scalar line, the slots being $u$'s undifferentiated scalar leg, $w$'s undifferentiated scalar leg and $w$'s second vector leg; \emph{BL}, a one-vector-leg vertex $u$ and a \cite[(1.6)]{Ba3} vertex $w$ joined by two scalar lines ending at legs of different index pairs of $w$, the slots being the two remaining legs of $w$ and $u$'s vector leg. In each shape every scalar slot is an undifferentiated leg and the vector slot carries no difference, as the degree $D=0$ requires. The bridged shapes are \cite[(2.21a)--(g)]{Ba3}; \cite[(2.21a)]{Ba3}, a vector bridge whose other side is a one-vector-leg vertex with a scalar loop, vanishes identically, and in the remaining shapes: a side of the bridge with no slot would be a one-slot subgraph, which vanishes (\cref{lem:a06:WNNC}(a)); a vector bridge is excluded: the side carrying the vector slot has two vector legs (the slot and the bridge), hence order $\ge2$, so the other side has order $1$ and is a single one-vector-leg vertex with no loop, whose two scalar legs are then both external, i.e.\ both scalar slots lie on it, and its difference acts on a slot, against $D=0$; so the bridge $l_0$ is scalar, each side carries one scalar slot, the side with the vector slot has order $\ge1$ and the other, a two-scalar-leg subgraph, order $\ge2$, whence orders $1+2$: the first side is a single one-vector-leg vertex $u$ whose difference, not allowed on its scalar slot or its vector leg, acts on $l_0$, and $S_2$ is one of the order-$2$ two-scalar-leg graphs, none of which has a bridge. This is the structure asserted in (c), and it is what Ba{\l}aban's arrows in \cite[(2.21b)--(g)]{Ba3} show.

\emph{(a): the common-anchor cancellation.} Let $\sigma_i\in u_i$ be the two scalar slots of $S$, on outer lines or external fields $F_i$. $G^{\mathrm{sw}}$ differs from $G$ only in the attachments at $\sigma_1,\sigma_2$, so $S$ has the same vertices, internal lines, differences, vertex localisations and anchor $\hat x$ in both, and the external lines of $S$ keep their indices, so that usefulness, the regrouped lines of the internal and external slices (which depend on the external indices only) and the transport of the multipliers to $\hat x$ (\cref{lem:a06d:S27}) agree. Hence the regrouped position kernel $K_S(y_1,y_2;\dots)$ of $S$, including its $R$-lines and the walled kernels, is the same function in both members, whatever the line kernels are. Its index structure is a matrix $M_S$ on the open index path from $\sigma_1$ to $\sigma_2$: every closed index cycle inside $S$ carries an even power of $q$ (else $\Tr q^{\mathrm{odd}}=0$ kills the term), and the total number of vector-leg ends at the vertices of $S$ is $2(\text{internal vector lines})+1$, odd, so the open path carries $q^{r}$ with $r$ odd times even polynomials in $q$, and $M_S^{\mathsf T}=-M_S$. At zeroth order the external fields are constants at $\hat x$ and the two local terms are $c(\hat x)F_1(\hat x)^{\mathsf T}M_SF_2(\hat x)$ and $c(\hat x)F_2(\hat x)^{\mathsf T}M_SF_1(\hat x)=-c(\hat x)F_1(\hat x)^{\mathsf T}M_SF_2(\hat x)$, with the same $c(\hat x)$; they cancel. Walled lines change nothing: they are scalar kernels times $\mathbf 1_N$ (\cref{lem:a06:WN1}(a)), so $M_S$ is unchanged; for slot fields which are kernels of several variables (the ends of $R^v_K$, $R^\gamma_K$), the zeroth-order term at $\hat x$ is their value at $\hat x$ in the slot variable, controlled by the norm \cite[(1.32)]{Ba3}. The remainders of the Taylor expansion are terms of the processing of Step~3(D$_4$), bounded term by term. This is Ba{\l}aban's reason for \cite[(2.18)]{Ba3}, ``the same expressions but with the opposite signs'' \cite[p.~444]{Ba3}, with the two members anchored at one vertex.

\emph{(b): the pinned root.} Here $u_1\ne u_2$ (the shapes T, B8, B10; in BL both scalar slots sit at $w$, so $u_1=u_2$ and (a) applies with the anchor at $w$). The local terms are $c_{S,1}(x)\Phi$ and $-c_{S,2}(x)\Phi$, with $c_{S,i}$ the regrouped coefficient with $u_i$ at $x$ and the anchor at $u_i$ (the sign from (a)). \emph{Split.} By the steps of \cref{thm:a06d:S48}: rescale to $\xi\Z^3$ ($\xi=L^{-j_0}$, $j_0$ the lowest external index), where by \cref{lem:a06:WN1}(f) and the rescaling of \cref{lem:a06:WN1}(a) every regrouped scalar line is $\breve G^v_l=C^\xi+R^v_l$ with a potential in $[0,\bar v]$ and $R^v_l$ as in \cref{lem:a06:WN1}(d), and every vector line is $C^\xi+R_l$ with the bulk remainder; move the multipliers to the anchor (\cref{lem:a06d:S27}); expand multilinearly, $c_{S,i}=c^C_{S,i}+c^R_{S,i}$. The shapes T, B8, B10 contain no \cite[(1.6)]{Ba3} or \cite[(1.7)]{Ba3} vertex, so no sharp indicator occurs and the internal sums are unrestricted on $\xi\Z^3$; and they contain no divergent proper subgraph (a two-scalar-leg or two-vector-leg subgraph inside a one-loop triangle or double line would need a bridge), so no inner counterterm occurs. \emph{The all-$C^\xi$ parts.} $c^C_{S,i}(x)$ is the same expression as in the bulk model: the same graph, vertex tensors and differences, the multipliers at the anchor, lines $C^\xi$, unrestricted sums. It is the zero-momentum value of one translation-invariant kernel on $\xi\Z^3$ with $u_i$ pinned: with $K^C_S(y_{u_1},y_{u_2},\dots)$ the all-$C^\xi$ kernel, $c^C_{S,1}(x)=\sum K^C_S(x,y_{u_2},\dots)=\sum K^C_S(0,y,\dots)$ and $c^C_{S,2}(x)=\sum K^C_S(y_{u_1},x,\dots)=\sum K^C_S(0,-y,\dots)$, both the total sum of $K^C_S(0,\cdot,\dots)$ over $\xi\Z^3$ (absolutely convergent at fixed $\xi$), so $c^C_{S,1}(x)=c^C_{S,2}(x)$, independent of $x$. The wall does not enter this term. \emph{The $R$-parts.} Each term of $c^R_{S,i}$ has at least one $R$-line. \Cref{lem:a06:WN1}(d) gives the bounds \eqref{eq:a06:WNR} for $R^v_l$ at every pair of points, with constant $C_R(1+\bar v)^3$, so $R^v_l$ is a remainder line of \cref{lem:a06d:S47}(b), and the step of \cref{thm:a06d:S48} for ``terms with at least one $R$-line, $D(S)=0$'' applies: the node is one-loop, its inner components are trees, the only dangerous ones are the \cite[(2.4)]{Ba3} subgraphs, treated by the summation by parts inside the coefficient (\cref{lem:a06d:S44}), every component after re-slicing has positive degree, and the bound function estimate \cite[(2.16)]{Ba3} bounds each $c^R_{S,i}$ by $C(1+\bar v)^p$ uniformly in $\xi$, $j_0$, $s$, $\eps$, the torus and $x$ (the kernel bounds are pointwise, and the anchor is fixed at $x$). The exponent $p$ counts the $R$-lines and the factors of $\bar v$ in their bounds, at most three lines here. \emph{Conclusion.} $|c_{S,1}(x)-c_{S,2}(x)|=|c^R_{S,1}(x)-c^R_{S,2}(x)|\le2C(1+\bar v)^p=: C^v_7$. This is the form of the $\phi\phi A$ clause of \cref{thm:a06d:S48}, $|c^{(0)}|\le C$, for the pair.

\emph{(c): bridged nodes.} Write $\Sigma_2(w,w')$ for the amputated kernel of $S_2$ between its leg point $w$ on the bridge and its other leg point $w'$, with $S_2$'s lines (walled or not, at the regrouped scale), and $F(y)\defeq(\partial^\xi_\mu\breve G_{l_0})(x_u,y)$ for the bridge factor, $u$'s difference acting on the bridge.

\emph{$u$ summed.} Choose the anchor $\hat x$ in $S_2$ (anchors of $\phi\phi A$ nodes are free, Step~3(D$_4$)). At zeroth order every external field of $S$ is a constant at $\hat x$, and so are $u$'s multipliers, which are smooth (\cref{lem:a06d:S27}; a one-vector-leg vertex carries no sharp indicator). The dependence of the regrouped kernel on $x_u$ is then the single factor $F$, and $\sum_{x_u}\xi^3(\partial^\xi_\mu\breve G_{l_0})(x_u,y)=0$ for every summable kernel: a forward difference in the summed variable telescopes, on $\xi\Z^3$ and on the torus. So $c=0$ for every member and any line kernels, before any $C^\xi/R$ split, and the same holds for the member with $S_2$ replaced by its counterterm vertex. No translation invariance of the lines is used.

\emph{$u$ pinned.} The anchor is $u$, at $x$. \emph{All-$C^\xi$ part.} On $\xi\Z^3$ with lines $C^\xi$ and unrestricted sums, $\sum_{w'}\Sigma^C_2(w,w')$ is a constant $g$ by translation invariance, and $\sum_w\xi^3(\partial^\xi_\mu C^\xi)(x,w)\,g=g\,\partial^\xi_\mu(C^\xi\mathbf 1)(x)=0$ since $C^\xi\mathbf 1\equiv1$. The counterterm member is the same computation with $\Sigma_2$ replaced by its own row sum. A sharp indicator inside $S_2$ occurs only at the \cite[(1.6)]{Ba3} tadpole, which cancels exactly against its counterterm member label by label ((Agr); the counterterm carries the same unit-cube label), so the pair has no indicator defect; the exchange bubble and the seagull loop carry only the smooth labels of their vector legs. So $c^C=0$ per member (and per pair). This is Ba{\l}aban's ``by translation invariance it can be written as a derivative of a constant'' \cite[p.~444]{Ba3} for \cite[(2.21c)]{Ba3}, and the exact cancellations he asserts for \cite[(2.21b), (2.21d), (2.21f), (2.21g)]{Ba3}, applied to the all-$C^\xi$ parts. \emph{Why the $R$-part is not bounded per member.} The term in which the bridge carries $R^v_{l_0}$ and $S_2$ carries lines $C^\xi$ equals $g\sum_w\xi^3(\partial^\xi_\mu R^v_{l_0})(x,w)=g\,\partial^\xi_\mu(R^v\mathbf 1)(x)$, where $g\sim\xi^{-1}$ (the node $S_2$ has $D(S_2)=-1$, and the $R$-line is not inside it, so it gains nothing for $S_2$). In the bulk $R\mathbf 1$ is constant ($-\Delta+m^2+aP$ maps constants to constants), and the term vanishes; under a wall $\breve G^v\mathbf 1=\mathbf 1-\breve G^v v$ is not constant, so the term is of order $\xi^{-1}$ and is not bounded uniformly in $\xi$. \emph{The pair.} Both members are in $\cC$ by (P). Let $f$ be the external field at $S_2$'s far slot $w'$ (constant at the anchor at zeroth order, Taylor-expanded at first order), and take the anchor convention (Agr) for $S_2$ with the anchor at the bridge end $w$, so that the counterterm member carries $\delta m^2_{S_2}(w)=\sum_{w'}\xi^3\Sigma_2(w,w')$, the row sum of the same kernel (the model's own counterterm, \cite[p.~417]{Ba3}). Then
\[
\sum_{w,w'}\xi^6F(w)\Sigma_2(w,w')f(w')-\sum_w\xi^3F(w)\,\delta m^2_{S_2}(w)f(w)\ =\ \sum_{w,w'}\xi^6F(w)\Sigma_2(w,w')\big[f(w')-f(w)\big],
\]
an identity for any kernels, with no symmetry used. Its zeroth-order term ($f$ constant) is $0$ exactly, including the terms in which an $R$-line sits inside $S_2$ (these have $D(S_2)+1=0$ and are paired with the counterterm member carrying the same $R$-line, as in the step of \cref{thm:a06d:S48} for $D(H)=-1$ components). The first-order term is $\sum_w\xi^3F(w)c_\mu(w)\,\partial_\mu f$ with $c_\mu(w)\defeq\sum_{w'}\xi^3\Sigma_2(w,w')(w'-w)_\mu$ the first moment of $S_2$; its all-$C^\xi$ part is bounded by \cite[(3.16)]{Ba3} (the cancellation $\sum_{y'}(\partial C^\xi\partial^*)(y'-y)=0$ of \cref{thm:a06d:S48}), and its $R^v$-part by $|\partial R^v\partial^*|\le C_R(1+\bar v)^3(\xi+\rho)^{-1}e^{-\delta_R\rho}$ (\cref{lem:a06:WN1}(d)) against the other line of the bubble, a sum convergent in $d=3$; so $|c_\mu(w)|\le C(1+\bar v)^3$, and $\sum_w\xi^3|F(w)|\le C(1+\bar v)^3$ by the bounds on $\partial C^\xi$ and $\partial R^v$. The H\"older remainder $f(w')-f(w)-(w'-w)\cdot\partial f(w)$ gives the node $S$ a differentiated external leg, after which every component has positive degree and \cite[(2.16)]{Ba3} bounds the sum. If the bridge attaches at $S_2$'s other leg and the anchor convention puts the counterterm at the far leg, the same identity holds with $f$ and $F$ exchanged, by the symmetry $\Sigma_2(w,w')=\Sigma_2(w',w)$ of the real symmetric scalar kernels at zero field (\cref{lem:a06:WN1}(a)). Altogether the $R$-part of the pair is $\le C(1+\bar v)^p=: C^v_{\mathrm{BR}}$, uniformly in $\xi$, $j_0$, $s$, $\eps$, the labels, the torus and $x$. The bridged shapes without a counterterm member (\cite[(2.21a), (2.21c), (2.21e)]{Ba3}: the one whose expression vanishes identically \cite[p.~442]{Ba3}, and those in which $S_2$ carries a difference on a leg or an extra power of $\eta$) need none: for them the all-$C^\xi$ part vanishes as above, and the $R$-part is convergent per member, since the difference on $S_2$'s leg, or the extra power of $\eta$, makes $D(S_2)\ge0$ before the gain of the $R$-line.
\end{proof}

\begin{remark}[reducible two-scalar-leg chains]\label{rem:a06:WNCH}
A one-particle-reducible two-scalar-leg chain of one-loop $D=-1$ parts is renormalised by the counterterms of its parts (Step~3(D$_4$)). The identity behind it, $T_{P_2}(b',c)=\delta_{b'c}\sum_{c'}K_{P_2}(b',c')$ with the same lines, indices and labels, holds for any kernels because the counterterm is computed with the model's own lines; so it holds verbatim in the walled model. The alternative route through translation invariance is not used.
\end{remark}
